\section{Introduction}
\label{sec:intro}

\emph{The claim.}
Paper~I~\cite{I} showed that an observer who looks at a periodic ripple through a lens of resolution $\sigma$ cannot beat a fixed exchange rate, wherever she puts her focus:
$g\sqrt\ell\ge1-\bar\varepsilon$, $\ell=\ln(1/\sigma)$.
It did so for windows whose ripple is weak in the sup norm, $\eta=\sup|R|\le1/\ell$.
The window class of Papers~F and~H~\cite{F,H} allows much stronger ripples: condition (P3), $\sum_nq0_n\le\Qb=\tfrac14$, gives only $\eta\le\tfrac14$.
This paper removes the assumption $\eta\le1/\ell$.
\begin{quote}
\emph{Wherever the observer focuses, whatever the strength of the ripple, the information that reaches her obeys $g\sqrt\ell\ge(1-\bar\varepsilon)(1-\eta)/(1+\eta)$.}
\end{quote}
The statement holds for every admissible window, every state and every $\ell\ge55.26$ (Theorem~\ref{thm:J}).
The factor $(1-\eta)/(1+\eta)$ is at least $3/5$, and it is one when the ripple vanishes.
It is what the ripple sandwich costs in this proof; whether the true quantity pays it is a separate question, which Sec.~\ref{sec:truth} answers numerically in the negative for the windows we could reach.

\emph{Where the factor comes from.}
A ripple of strength $\eta$ multiplies the window by a factor between $1-\eta$ and $1+\eta$.
Paper~H~\cite{H} turned this into a sandwich: the gains of the expansion are at least $\lambda=1/(1+\eta)$ times their ripple-free values, the cost in the denominator at most $\Upsilon=1/(1-\eta)$ times.
So the natural law for a strong ripple has the coefficient $\lambda/\Upsilon=(1-\eta)/(1+\eta)$.
Paper~H proved it for the states whose focus spans at least about $0.7$ ripple periods and $\ell\le960.8$ (Theorem~H1, (ii)--(iii)).
Paper~I proved the law with coefficient one, on the whole map of states, for a weak ripple.
Here the two meet: the whole map, every ripple of the class, coefficient $\lambda/\Upsilon$.
In the lower bounds of this proof the coefficient cannot be one at $\eta=\tfrac14$: with the target unscaled their margin falls from $+24.5$ to $-26.5$ at $\ell=55.26$ \numst.
The sandwich puts the worst gain and the worst cost of the ripple at the same time, and one window need not have both.
The true quantity, computed on windows with $\eta=\tfrac14$, stays far above the law ($((a)+(a'))\sqrt\ell\ge3.40$ on everything we found at $\ell=55.26$, so $g\sqrt\ell\ge(1-\bar\varepsilon)\,3.40$ there), and at those states the factors of the sandwich move it by a few per cent (Sec.~\ref{sec:truth}) \numst.
We conjecture that the coefficient one holds for every ripple of the class (Conjecture~\ref{conj:one}).

\emph{What is new.}
Most of the work of Paper~I carries over.
Its lines (the lower bounds of the local expansion) are written with $\eta$ as a parameter.
We read every lemma behind them again, those of Paper~I and those it borrows from Papers~F and~H, and recorded for each where $\eta$ enters its proof and which bound on $\eta$ that step uses (Table~\ref{tab:lemmas}, App.~\ref{app:lemmas}): every lemma uses only $\eta<1$ (through $h\ge1-\eta$) or no bound at all, with one exception.
What had to be added is the following.
\begin{itemize}
\item \emph{Two ends instead of a continuum} (Sec.~\ref{sec:sandwich}).
  The remainders of the lines are functions of $\eta$; the gains and the target carry the sandwich.
  Each remainder is a sum of products of six non-negative, increasing, convex functions of $\eta$ (Lemma~\ref{lem:chord}); as a function of $\lambda=1/(1+\eta)$ it is convex and lies under its chord on $[\tfrac45,1]$.
  With the target scaled by $\lambda/\Upsilon=2\lambda-1$, every line becomes affine in $\lambda$, and the interval proofs at the two ends, $\eta=0$ and $\eta=\tfrac14$, cover every ripple in between.
\item \emph{One step repaired} (Sec.~\ref{sec:changes}).
  One elementary step of Paper~I (in the group FrB1) used $\eta\le1/\ell$; it fails at $\eta=\tfrac14$.
  A larger constant $X_1\to U_1/(1-\bar\eta)+\sqrt{\mathrm{CZ}(\bar\eta)}$ repairs it; no new lemma is needed.
\item \emph{New constants} (Sec.~\ref{sec:changes} and App.~\ref{app:D}).
  The rule of the proof, which chooses the option of the expansion, has new constants ($c_1=0.77$, $c_T=6.5$, and $c_{\rm ch}=1.5$ for $\ell<960.8$, $2.0$ beyond), $18$ of the $48$ quoted bounds of Paper~I take new values, and the tail $[10^6,\infty)$ is derived again with $\eta$ and $1/\ell$ separated.
\end{itemize}
The focus at the resolution keeps the coefficient one (Theorem~\ref{thm:Jres}): $(a)+(a')\ge0.49$ whenever the ripple reaches the focus ($G>0$), and $\ge0.1995$ otherwise, for every ripple of the class; both are above $\Gamma_\ell\le0.1462$, by hand \proofst.
The very wide focus needed no bound on $\eta$ already in Paper~I.

\emph{What is proved and what is not.}
We keep the distinction of Papers~F, H and~I between proof, computer-assisted proof, derivation, numerical evidence and conjecture, and mark every statement.
The computer-assisted part is a finite list of interval-arithmetic evaluations with outward rounding, at the two ends; it is listed in App.~\ref{app:D}.
We tried to break the claim first, with planted dips and with deliberately wrong versions that the same search breaks (Sec.~\ref{sec:break}).
The limits are those of Paper~I with one removed (Sec.~\ref{sec:limits}): the class constant $\bar\varepsilon$, the upper side, and now also the restriction of the coefficient to the ripple sandwich, since Paper~H's general sandwich factors do not close the proof here.

\emph{In the observer's words.}
A stronger ripple can dim the picture at most by the ratio $(1-\eta)/(1+\eta)$ of the darkest to the brightest it can make it.
It does not change the exchange rate.
However the observer focuses, the information that survives the blur is still at least that fraction of $1/\sqrt\ell$; where she focuses at the resolution itself, the ripple cannot even take that fraction away; and in the windows we could compute, it does not take it away anywhere.
